\begin{table}[t]
\centering
\small
\setlength{\tabcolsep}{6pt}
\begin{tabular}{lccc|ccc|ccc}
\toprule
\multirow{2}{*}{\shortstack{\textbf{Eval.}\\\textbf{Modality}}}
 & \multicolumn{3}{c|}{\textbf{Training Modality}}
 & \multicolumn{3}{c|}{\textbf{VisDoc Training Source}}
 & \multicolumn{3}{c}{\textbf{Video Training Source}} \\
\cmidrule(lr){2-4} \cmidrule(lr){5-7} \cmidrule(lr){8-10}
 & Image & VisDoc & Video
 & ColPali & VisRAG$_{\text{ind}}$ & VisRAG$_{\text{syn}}$
 & Vid2Cap & Cap2Vid & VideoQA \\
\midrule
Image   & 58.45 & 39.87 & 30.64
        & 32.29 & 4.14  & 3.53
        & 20.48 & 3.52  & 4.12 \\

VisDoc  & 58.73 & 61.72 & 41.55
        & 53.66 & 1.26  & 0.79
        & 34.23 & 0.80  & 0.56 \\

Video   & 33.58 & 23.62 & 26.30
        & 25.40 & 10.86 & 10.76
        & 23.89 & 10.43 & 10.23 \\
\midrule
AVG     & 52.80 & 39.87 & 30.64
        & 37.27 & 4.80  & 4.36
        & 25.50 & 4.28  & 4.43 \\
\bottomrule
\end{tabular}
\caption{
Performance comparison across evaluation modalities (rows) under different training data settings (columns).Rows indicate the evaluation modality of the test datasets (Image, VisDoc, or Video).
The first column block compares models trained with a single modality to analyze modality-level training effects. The second block varies VisDoc training data sources to study the quality of various VisDoc data sources. The third block varies Video training data sources to analyze the impact of video data sources. For fair comparison, each training configuration uses 100K training samples.
}
\label{tab:data_mix_t1}
\end{table}
